\lecture{11}
\title{The Recursion Theorem}

\begin{document}

\subsection{Quines: Python}

First, we construct a Turing Machine \textbf{quine}---a
Turing Machine that takes nothing as input
and prints out its own description.

As a warmup, here's an example of a quine in Python.

\begin{verbatim}[python]
lines = [
    'lines = [',
    ']',
    'quote = chr(39)',
    'print(lines[0])',
    'for line in lines:',
    '    print("    " + quote + line + quote + ",")',
    'for line in lines[1:]:',
    '    print(line)',
]
quote = chr(39)
print(lines[0])
for line in lines:
    print("    " + quote + line + quote + ",")
for line in lines[1:]:
    print(line)
\end{verbatim}

Note that \texttt{chr(39)} on line 11
is just a Python-safe way of writing a single quote.
Here's how this quine works.
\begin{itemize}
    \item \textbf{Part A: Lines 1–10.}
    When Part A is run,
    it writes a description $D$
    to a list \texttt{lines}.
    Note that $D$ consists of the following:
    \begin{itemize}
        \item (Line 2) First, a header \texttt{lines = [}.
        \item (Line 3) Then, a footer \texttt{]}.
        \item (Lines 4–9) Finally, a description of Part B.
    \end{itemize}
    \item \textbf{Part B: Lines 11–16.}
    When Part B is run, it does the following:
    \begin{enumerate}
        \item First, it reads the description $D$
        that has been written to \texttt{lines}.
        \item (Lines 12–14) Then, it prints lines of Python code
        that would print $D$; that is, it prints Part A.
        \item (Lines 15–16) Finally, it prints the footer and body (i.e., Lines 3–9) of $D$.
    \end{enumerate}
\end{itemize}

We'll follow a similar structure
when constructing a quine in the language of Turing Machines.

\subsection{Quines: Turing Machines}

\textbf{Definition.}
We require the following two definitions.
\begin{itemize}
    \item For any string $w$,
    let $P_w$ denote the Turing machine that, on any input,
    writes $w$ to the tape, then halts.
    \item Let $q \colon \Sigma^* \to \Sigma^*$ be the (computable) function
    defined by $q(w) := \langle P_w \rangle$.
\end{itemize}

\textbf{Example.}
In the Python analogy, we have:
\begin{itemize}
    \item Part A is just $P_D$;
    when Part A is run, it writes the string $D$
    to \texttt{lines} (or, equivalently, the tape).
    \item After reading in $D$, lines 12–14 of Part B
    simply compute $q(D)$ and print it.
    \begin{itemize}
        \item Note that $q(D)$ is just a description of Part A.
    \end{itemize}
\end{itemize}

We're now ready to describe our construction of
a Turing Machine quine!
(We denote this construction by $\mathrm{SELF}$.)

\begin{quote}
    \textbf{Construction ($\mathbf{SELF}$).}
    The Turing machine $\mathrm{SELF}$ consists of two parts
    $A$ and $B$.
    
    Define $A := P_{\langle B \rangle}$,
    and define $B$ as follows:
    \begin{enumerate}
        \item First, it reads the description $\langle M \rangle$
        of a portion of a TM that has been written to the tape.
        \item Then, it prints a description of a Turing machine
        that would write $\langle M \rangle$ to the tape;
        that is, it prints $q(\langle M \rangle)$.
        \item Finally, it prints $\langle M \rangle$.
    \end{enumerate}
    Now notice that since $A$ was run before $B$,
    the input $\langle M \rangle$ that $B$ reads off of the tape
    is exactly $\langle B \rangle$.

    Therefore, what $B$ prints is $q(\langle B \rangle)
    = \langle P_{\langle B \rangle} \rangle = \langle A \rangle$
    followed by $\langle B \rangle$,
    which is just $\langle \mathrm{SELF} \rangle$.  Done!
\end{quote}

Notice the direct parallels with the Python quine.
The only major difference is that
the Turing machine quine doesn't need to include
the header \texttt{lines = [} and the footer \texttt{]}.

\begin{remark}
    Technically, what $B$ should really print out is
    the description of the Turing machine given by
    the combination of $q(\langle M \rangle)$ and $\langle M \rangle$.
    This is necessary in case the description of
    the combination of two Turing machines
    is not the same as the combination of their descriptions.
\end{remark}


For an English-language example of a quine
that follows a similar structure, consider the following sentence.
\begin{quote}
    Print two copies of the following,
    the second one in quotes:
    ``Print two copies of the following,
    the second one in quotes:''
\end{quote}
Note that the following, in contrast, is \emph{not} a quine;
its usage of ``this sentence'' is circular.
\begin{quote}
    Print out this sentence.
\end{quote}

\subsection{The Recursion Theorem:
$A_{\mathrm{TM}}$, $\mathrm{MIN}_{\mathrm{TM}}$}

We can apply this self-reference trick
more generally in a lot more settings than just printing.

\begin{theorem}{The Recursion Theorem}
    Let $T$ be a Turing machine
    that computes a function
    $t \colon \Sigma^* \times \Sigma^* \to \Sigma^*$.
    Then there is a Turing machine $R$
    that computes a function $r \colon \Sigma^* \to \Sigma^*$
    where, for every $w \in \Sigma^*$, we have
    $r(w) = t(\langle R \rangle, w)$.
    \begin{quote}
        \emph{Informal Rephrasing.}
        Any TM $R$ is allowed to
        acquire its own description $\langle R \rangle$
        and perform computations with it.
    \end{quote}
\end{theorem}

\begin{proof}
    Construct $R$ in three parts $A$, $B$, and $T$
    (where $T$ is the TM given in the theorem statement).
    Let's redefine some notation:
    \begin{itemize}
        \item Let $P_w$ denote the Turing machine that
        writes $w$ on the tape
        \emph{after what is already written}, then halts.
        \item Define $q \colon \Sigma^* \to \Sigma^*$
        by $q(w) := \langle P_w \rangle$, just as before.
    \end{itemize}
    We then define $A := P_{\langle BT \rangle}$,
    and define $B$ as follows:
    \begin{enumerate}
        \item First, it reads the description $\langle M \rangle$
        currently written to the tape.
        \item Then, it computes $q(\langle M \rangle)$
        and combines it with $\langle M \rangle$
        to obtain a TM description $\langle X \rangle$.
        \item Finally, it writes $\langle X, w \rangle$ to the tape.
    \end{enumerate}
    As before, when $B$ is run,
    it will read $\langle M \rangle = \langle BT \rangle$, and so
    it will write $\langle ABT, w \rangle
    = \langle R, w \rangle$ to the tape.

    Thus, when $T$ runs after $B$,
    it will have both $w$ and the description $\langle R \rangle$
    to work with as input, as desired. ~ $\blacksquare$
\end{proof}

The recursion theorem essentially packages together
all of the self-referential work we did earlier
into a reusable tool.

\begin{problem}
    Describe a Turing machine $\mathrm{SELF}$
    that prints its own description.
\end{problem}

\begin{solution}
    Here is the description of $\mathrm{SELF}$.
    \begin{enumerate}
        \item Obtain, via the recursion theorem,
        the self-description $\langle \mathrm{SELF} \rangle$.
        \item Print $\langle \mathrm{SELF} \rangle$.
    \end{enumerate}
    Explicitly, we are applying the recursion theorem when
    $T$ is the machine that, on input $\langle M, w \rangle$,
    just prints $\langle M \rangle$. ~ $\blacksquare$
\end{solution}

Now for some less immediate examples.

\begin{theorem}{$A_{\mathrm{TM}}$ Undecidable}
    The language $A_{\mathrm{TM}} := \{\langle M, w \rangle \mid
    M \text{ is a TM that accepts } w\}$ is undecidable.
\end{theorem}

\begin{proof}
    Suppose there is a decider $H$ for $A_{\mathrm{TM}}$.
    Then consider the TM $B$ that, on input $w$,
    does the following:
    \begin{enumerate}
        \item Obtain, via the recursion theorem,
        the self-description $\langle B \rangle$.
        \item Run $H$ on input $\langle B, w \rangle$.
        If $H$ accepts, then reject.  If $H$ rejects, then accept.
    \end{enumerate}
    Now running $B$ on any input $w$
    raises a contradiction. ~ $\blacksquare$
\end{proof}

\begin{remark}
    Compare with the diagonalization argument used
    back in \href{/18.404/lectures/8/#a-tm-undecidability}{Lecture 8}
    to prove the same fact.
\end{remark}

\begin{theorem}{$\mathrm{MIN}_{\mathrm{TM}}$ Unrecognizable}
    We first define the notion of \textbf{minimal} Turing machines.
    \begin{quote}
        \textbf{Definition.}
        Say a TM $M$ is \textbf{minimal} if
        there is no TM $M'$ equivalent to $M$
        such that $\langle M' \rangle$
        has shorter length than $\langle M \rangle$.
    \end{quote}
    Then the language $\mathrm{MIN}_{\mathrm{TM}} := \{ \langle M \rangle \mid
    M \text{ is a minimal TM}\}$ is not Turing-recognizable.
\end{theorem}

\begin{proof}
    Recall from \href{/18.404/lectures/6/#enumerators-turing-machines-with-printers}{Lecture 6}
    that $\mathrm{MIN}_{\mathrm{TM}}$ is Turing-recognizable
    if and only if some enumerator $E$ prints out $\mathrm{MIN}_{\mathrm{TM}}$.

    Suppose (FTSOC) that such an enumerator $E$ exists.
    Then consider the TM $C$ that, on input $w$,
    does the following:
    \begin{enumerate}
        \item Obtain, via the recursion theorem,
        the self-description $\langle C \rangle$.
        \item Run the enumerator $E$ until
        a machine $D$ appears with a longer description
        than that of $C$.
        \item Simulate $D$ on input $w$.
    \end{enumerate}
    Then $C$ is equivalent to $D$ and is shorter than it,
    so $D$ is not minimal.
    But $E$ printed out $D$, contradiction. ~ $\blacksquare$
\end{proof}

One more example, of a slightly different flavor.

\begin{theorem}{Fixed Points}
    For any computable function $t \colon \Sigma^* \to \Sigma^*$,
    there exists a TM $F$ for which
    $t(\langle F \rangle)$ describes a TM equivalent to $F$.
\end{theorem}

\begin{proof}
    Define $F$ so that, on input $w$, it does the following:
    \begin{enumerate}
        \item Obtain, via the recursion theorem,
        the self-description $\langle F \rangle$.
        \item Compute $t(\langle F \rangle)$
        to obtain the description of a TM $G$.
        \item Simulate $G$ on $w$.
    \end{enumerate}
    Then clearly $\langle F \rangle$ and
    $t(\langle F \rangle) = \langle G \rangle$
    describe equivalent Turing machines,
    because $F$ simulates $G$. ~ $\blacksquare$
\end{proof}

\end{document}